\documentclass[10pt,a4paper]{amsart}
\usepackage[margin=19mm]{geometry}
\usepackage{amsmath,amssymb,mathtools}
\usepackage[scr=rsfs,cal=euler]{mathalfa}
\usepackage{xfrac}
\usepackage[unicode,hidelinks,bookmarks=false]{hyperref}
\newcommand{\ie}{i.e.\ }
\newcommand{\cf}{cf.\ }
\newcommand{\resp}{respectively\ }
\newcommand{\Subsection}{\subsection}
\providecommand{\nobreakdash}{\nobreak\hbox{-}}
\setlength{\emergencystretch}{2em}
\multlinegap0pt
\usepackage{bm}
\usepackage{bbm}
\usepackage{dsfont}
\usepackage{xspace}
\usepackage{cancel}
\usepackage{mathtools}
\usepackage{amsthm}
\makeatletter
\@ifundefined{lem}{\newtheorem{lem}{Lemma}}{}
\makeatother
\title[Degree-truncated choosability of graphs]{Degree-truncated choosability of graphs}
\author{Huan Zhou, Jialu Zhu, Xuding Zhu}
\date{Source: arXiv:2507.10453v1 (CC BY 4.0)}
\begin{document}
\maketitle

\noindent\emph{Source and licence.} Degree-truncated choosability of graphs. Version arXiv:2507.10453v1 (CC BY 4.0), distributed under the Creative Commons Attribution 4.0 International licence, as recorded in the arXiv licence metadata for this exact version and shown on the abstract page https://arxiv.org/abs/2507.10453v1. Full rights notice: licenses/arxiv-2507.10453-CC-BY-4.0.txt.\par\smallskip

\noindent\textit{Editorial scope.} Counted unit: the Lem whose source label is \texttt{lem-verynice} in the article (source0.tex lines 442--445, source label \texttt{lem-verynice}), delivered together with the article's own complete original proof of it (source0.tex lines 446--504, one \texttt{proof} environment of 59 lines). No mathematics is written, repaired or re-derived: statement and proof are the article's own text line for line. Required context retained uncounted: none. Every definition, display and label of those lines is retained unchanged. Omitted from this package and not redistributed: 1--441, 505--748. Nothing inside a retained excerpt is omitted or altered: every retained body is delivered exactly as the article's lines are written, so no omission receipt of any kind accompanies this package. Citations: the retained excerpts carry no citation command at all, so no bibliography entry of the article is transcribed and no reference is redistributed. Reference adaptation: none was needed and none was made; every reference inside the delivered unit resolves inside it, so no unresolved reference or citation is produced. Printed slips kept literal: no printed word, symbol, hypothesis or number of the counted statement or of its proof was altered. Layout and class: the article's own class is \texttt{article} with packages including mathrsfs, amsmath, amsthm, amssymb, bm, mathtools, thm-restate, ulem, cancel, geometry, fontenc, enumitem, natbib, todonotes; the wrapper is the lane's pinned \texttt{amsart} editorial wrapper at 10pt on A4, which restates the theorem-like environments the retained text needs and adds the article's own macro definitions that the retained text needs (none), with extra wrapper packages (bm, bbm, dsfont, xspace, cancel, mathtools). The delivered numbering is the unit's own. Assets: no retained span contains a figure environment, a graphics-inclusion command, a PDF-page inclusion, a table or a TikZ picture; no image file is copied into this package, the package declares no asset dependency and no image file is redistributed with it. Licence: version arXiv:2507.10453v1 of this article is distributed under the Creative Commons Attribution 4.0 International licence, as recorded in the arXiv licence metadata for this exact version and shown on the abstract page https://arxiv.org/abs/2507.10453v1. A full rights notice is delivered at licenses/arxiv-2507.10453-CC-BY-4.0.txt. Characters: every retained excerpt is pure ASCII, so the delivered engine, XeLaTeX with its default Latin Modern text font, reports no missing character.
\begin{lem}
    \label{lem-verynice}
    For every  plane graph $G$, for any vertex $v^*$ on the boundary of infinite face of $G$, $\Theta(G)$ has a very nice subgraph with respect to $v^*$.
\end{lem}

\begin{proof}
   We shall prove Lemma \ref{lem-verynice} by induction on the number of vertices of $G$. 
   

   \medskip
   \noindent
   {\bf Case 1} $G$ is 2-connected.
   \medskip
   
   If $G=K_1$ or $K_2$, then $H=\Theta(G)$ is very nice with respect to $v^*$. 
   
   If $G \ne K_1, K_2$ is an outerplanar graph, then $G$ contains a cycle $C=(v_1,v_2,\ldots, v_k)$, bounding a finite face $\theta$,  such that $d_G(v_i)=2$ for $i=2,3,\ldots, k-1$, and $v^* \notin \{v_2, v_3, \ldots, v_{k-1}\}$. Let $G'=G-\{v_2, v_3, \ldots, v_{k-1}\}$. Then $\Theta(G)$ is obtained from $\Theta(G')$ by adding vertices $v_2,v_3, \ldots, v_{k-1}$ and $\theta$, and edges $v_i \theta$ for $i=1, 2, \ldots, k$ and $v_i \theta^*_G$ for $i=2,3,\ldots, k-1$. 
   
   By the induction hypothesis, $\Theta(G')$ has a very nice subgraph $H'$  with respect to $v^*$. Let $H$ be obtained from $H'$ by adding vertices $v_2,v_3, \ldots, v_{k-1}$ and $\theta$, and edges $v_i \theta, v_i \theta^*_G$   for $i=2,3,\ldots, k-1$.  Then $H$ is a very nice subgraph of $\Theta(G)$ with respect to $v^*$. 



   Assume $G$ is not an outerplanar graph. 
   If $G$ has a degree 2 vertex $v \ne v^*$ not contained in a triangle, then let $G'$ be obtained from $G$ by removing the vertex $v$ and adding an edge connecting the two neighbours of $v$. Then $F(G)=F(G')$ and $\Theta(G)$ is obtained from $\Theta(G') $ by adding the vertex $v$ and the edges $v \theta_1, v \theta_2$, where $\theta_1,\theta_2$ are the two faces of $G$ incident to $v$. By induction hypothesis, $\Theta(G')$ has a very nice subgraph $H'$ with respect to $v^*$. Let $H$ be obtained from $H'$ by adding vertex $v$ and edges $v \theta_1, v \theta_2$. Then $H$ is a very nice subgraph of $\Theta(G)$ with respect to $v^*$.
 
Otherwise let $u$ be a vertex not on the boundary cycle of $G$, and let $G'=G-u$. Let $\theta_u$ be the face of $G'$ containing the vertex $u$. 
Let $C$ be the boundary cycle of $\theta_u$, and let 
 $u_1,u_2,\ldots, u_k$ be the neighbours of $u$, occurring in $C$ in this cyclic order in clockwise direction.  For $i=1,2,\ldots, k$, let $P_i$ be the subpath of $C$ from $u_i$ to $u_{i+1}$ (where let $u_{k+1}=u_1$) along the clockwise direction. So $u_i,u_{i+1}$ are the end vertices of $P_i$, and $\cup_{i=1}^k(P_i-\{u_i\})$ is a partition of $V(C)$. 
    
    
    In $G$, $\theta_u$ is divided into $k$ faces $\theta_1, \theta_2, \ldots, \theta_k$, where $V(\theta_i) = V(P_i) \cup \{u\}$. 
    
    Let $H'$ be a very nice subgraph of $\Theta(G')$ with respect to $v^*$. Let $Z=N_{\Theta(G')}(\theta_u) -N_{H'}(\theta_u)$. Then $|Z| \le 2$. 
    
    Assume $Z \subseteq \{z_1,z_2\}$. Assume   $z_1 \in P_i - \{u_i\}$, and $z_2 \in P_j-\{u_j\}$. 
    If $i \ne j$, then let $H$ be obtained from $H'$ by deleting  $\theta_u$ (and hence all the  edges $\{v\theta_u: v \in V(\theta_u)\}$), and adding vertices $u, \theta_1,\theta_2,
    \ldots, \theta_k$ and adding  edges 
    $$\cup_{t=1}^k\{v\theta_t: v \in V(P_t)-\{u_t\}\} \cup \{u\theta_i,u\theta_j\} - \{z_1\theta_i,z_2\theta_j\}.$$ 
    If $i=j$, then without loss of generality, assume that 
$i=j=1$. Let 
$H$ be obtained from $H'$ by deleting  $\theta_u$ (and hence all the  edges $\{v\theta_u: v \in V(\theta_u)\}$), and adding vertices $u, \theta_1,\theta_2,
    \ldots, \theta_k$ and adding  edges   
    $$\cup_{t=1}^k\{v\theta_t: v \in V(P_t)-\{u_t\}\} \cup \{u\theta_1, u\theta_k, u_1\theta_1\} - \{z_1\theta_1,z_2\theta_1, u_1\theta_k\}.$$

    It is straightforward to verify that $H$ is a very nice subgraph of $\Theta(G)$ with respect to $v^*$.  


\medskip
\noindent
{\bf Case 2} 
 $G$ is not 2-connected.
 \medskip

Note that $G$ may be connected or not connected. Let $B$ be a leaf block of $G$ (which could be a 2-connected component of $G$) such that
all vertices of $G-U(B)$ are contained in the infinite face of $B$ and $v^* \notin V(B)$. It is obvious that such a leaf block $B$ exists. Let $G'=G-U(B)$. By the induction hypothesis, $\Theta(G')$ has a very nice subgraph $H'$ with respect to $v^*$.  

Let $\theta_B$ be the face of $G'$ that contains $U(B)$. 

If $B$ is a 2-connected component of $G$, then $\Theta(G)$ is obtained from the disjoint union of $\Theta(G')$ and $\Theta(B)$ by identifying $\theta_B$ with $\theta^*_B$. Let $H''$ be a very nice subgraph of $\Theta(B)$ with respect to $v'$ for an arbitrary boundary vertex $v'$ of $B$. Let $H$ be obtained from the disjoint union of $H'$ and $H''$ by identifying $\theta_B$ with $\theta^*_B$. Then $H$ is a very nice subgraph of $\Theta(G)$ with respect to $v^*$. 

Otherwise $B$ contains a cut-vertex $v'$ of $G$, and $\Theta(G)$ is obtained from the disjoint union of $\Theta(G')$ and $\Theta(B)$ by identifying $\theta_B$ with $\theta^*_B$, and identifying the copy of $v'$ in $G'$ and the copy of $v'$ in $B$. Let $H''$ be a very nice subgraph of $\Theta(B)$ with respect to $v'$. Let $H$ be obtained from the disjoint union of $H'$ and $H''$ by identifying $\theta_B$ with $\theta^*_B$, and identifying the copy of $v'$ in $H'$ and the copy of $v'$ in $H''$.   Then $H$ is a very nice subgraph of $\Theta(G)$ with respect to $v^*$. 

This completes the proof of Lemma \ref{lem-verynice}.
\end{proof}

\end{document}
